\documentclass[11pt]{article}
\usepackage[margin=1in]{geometry}
\usepackage{amsmath,amssymb,amsthm}
\usepackage{array}
\usepackage{parskip}

\newtheorem{theorem}{Theorem}
\newtheorem{lemma}[theorem]{Lemma}
\newtheorem{proposition}[theorem]{Proposition}
\newtheorem{corollary}[theorem]{Corollary}
\theoremstyle{definition}
\newtheorem{definition}[theorem]{Definition}
\theoremstyle{remark}
\newtheorem{remark}[theorem]{Remark}

\title{Disproof of conjecture \texttt{00000000443}}
\author{earthking11}
\date{2026-09-12}

\begin{document}
\maketitle

\section*{Verdict: FALSE}

We disprove conjecture \texttt{00000000443} as stated. There are two
\emph{independent} failures:

\begin{itemize}
\item[(A)] the asserted lower bound $\dim L_n \ge 2^{n-1} - 2^{\lceil n/2\rceil}$
  is false for \emph{every} $n \ge 4$ (it first fails at $n = 4$);
\item[(B)] the asserted prime-gap formula
  ``the gap is exactly half of $2^{(n-1)/2}$'' is false at $n = 3, 7, 11$
  (it happens to hold at $n = 5$, and is not even integral at $n = 2$).
\end{itemize}

Each failure alone falsifies the conjecture. The exact dimensions used below are
the standard ones, given by Witt's formula; the disproof is pure arithmetic and
requires no structure theory.

\section{The conjecture}

\begin{quote}
\textbf{English.} Definition: The $n$-th homogeneous dimension of the free Lie
algebra $L(V)$ is given by Witt's formula
$\frac1n\sum_{d\mid n}\mu(d)k^{n/d}$. Conjecture: For $k = 2$ and all $n \ge 2$,
$\dim L_n \ge 2^{n-1} - 2^{\lceil n/2\rceil}$, and for $n$ prime the gap between
this bound and the exact value is exactly half of $2^{(n-1)/2}$.
\end{quote}

The statement as filed in \texttt{conjectures/00000000443.md} also carries a Chinese
translation; it is reproduced in the accompanying \texttt{README.md}.

Throughout, $L_n$ denotes the degree-$n$ homogeneous component of the free Lie
algebra on two generators ($k = 2$), and $\mu$ is the Möbius function.

\section{Witt's formula and the exact dimensions}

\begin{theorem}[Witt]
\label{thm:witt}
For the free Lie algebra on $k$ generators, the homogeneous component of degree
$n \ge 1$ has dimension
\[
\dim L_n \;=\; \frac1n \sum_{d \mid n} \mu(d)\, k^{\,n/d}.
\]
\end{theorem}

We use the specialisation $k = 2$. The values for $n = 1,\dots,12$ are given in
Table~\ref{tab:dim}. For example, at $n = 4$ the divisors are $1,2,4$ with
$\mu(1)=1$, $\mu(2)=-1$, $\mu(4)=0$, so
\[
\dim L_4
= \frac14\bigl(\mu(1)2^{4} + \mu(2)2^{2} + \mu(4)2^{1}\bigr)
= \frac14(16 - 4 + 0)
= \frac{12}{4}
= 3 .
\]

\begin{table}[h]
\centering
\begin{tabular}{r|r|r|r|r|r|r|r|r|r|r|r|r}
$n$ & 1 & 2 & 3 & 4 & 5 & 6 & 7 & 8 & 9 & 10 & 11 & 12 \\
\hline
$\dim L_n$ & 2 & 1 & 2 & 3 & 6 & 9 & 18 & 30 & 56 & 99 & 186 & 335
\end{tabular}
\caption{Exact dimensions $\dim L_n$ for $k=2$, $n = 1,\dots,12$, from
Witt's formula (Theorem~\ref{thm:witt}).}
\label{tab:dim}
\end{table}

\section{The claimed bound}

Write
\[
B(n) \;=\; 2^{\,n-1} - 2^{\lceil n/2\rceil},
\]
so the conjecture asserts $\dim L_n \ge B(n)$ for all $n \ge 2$.
Table~\ref{tab:bound} compares $B(n)$ with $\dim L_n$.

\begin{table}[h]
\centering
\begin{tabular}{r|r|r|c}
$n$ & $\dim L_n$ & $B(n) = 2^{n-1} - 2^{\lceil n/2\rceil}$ & $\dim L_n \ge B(n)$? \\
\hline
1  & 2   & $2^{0}-2^{1} = -1$                    & yes \\
2  & 1   & $2^{1}-2^{1} = 0$                     & yes \\
3  & 2   & $2^{2}-2^{2} = 0$                     & yes \\
4  & 3   & $2^{3}-2^{2} = 4$                     & \textbf{no} \\
5  & 6   & $2^{4}-2^{3} = 8$                     & \textbf{no} \\
6  & 9   & $2^{5}-2^{3} = 24$                    & \textbf{no} \\
7  & 18  & $2^{6}-2^{4} = 48$                    & \textbf{no} \\
8  & 30  & $2^{7}-2^{4} = 112$                   & \textbf{no} \\
9  & 56  & $2^{8}-2^{5} = 224$                   & \textbf{no} \\
10 & 99  & $2^{9}-2^{5} = 480$                   & \textbf{no} \\
11 & 186 & $2^{10}-2^{6} = 960$                  & \textbf{no} \\
12 & 335 & $2^{11}-2^{6} = 1984$                 & \textbf{no} \\
\end{tabular}
\caption{The claimed bound $B(n)$ against the exact dimension. The inequality
holds only for $n = 2, 3$ and fails for every $4 \le n \le 12$. (At $n=1$ the
formula gives a negative value, but the conjecture only ranges over $n \ge 2$.)}
\label{tab:bound}
\end{table}

\section{Failure (A): the inequality is false for every $n \ge 4$}

\begin{lemma}
\label{lem:dim4}
$\dim L_4 = 3$.
\end{lemma}

\begin{proof}
By Theorem~\ref{thm:witt} with $k = 2$ and $n = 4$: the positive divisors of $4$
are $1,2,4$, and $\mu(1)=1$, $\mu(2)=-1$, $\mu(4)=0$. Hence
\[
\dim L_4 = \frac14\bigl(1\cdot 2^{4} + (-1)\cdot 2^{2} + 0\cdot 2^{1}\bigr)
= \frac14(16-4) = \frac{12}{4} = 3. \qedhere
\]
\end{proof}

\begin{lemma}
\label{lem:b4}
$2^{\,3} - 2^{\,2} = 4$, i.e.\ $B(4) = 4$.
\end{lemma}

\begin{proof}
$2^{3} = 8$ and $2^{2} = 4$, so $2^{3} - 2^{2} = 8 - 4 = 4$. Since
$\lceil 4/2 \rceil = 2$, this is exactly $B(4)$.
\end{proof}

\begin{theorem}[failure of the inequality]
\label{thm:Afails}
$\dim L_4 = 3 < 4 = B(4)$. Consequently the conjecture's assertion
``$\dim L_n \ge 2^{n-1} - 2^{\lceil n/2\rceil}$ for all $n \ge 2$'' is false.
\end{theorem}

\begin{proof}
By Lemma~\ref{lem:dim4}, $\dim L_4 = 3$. By Lemma~\ref{lem:b4}, $B(4) = 4$.
Since $3 < 4$, the inequality fails at $n = 4$. A single counterexample refutes a
universal statement. In fact Table~\ref{tab:bound} shows the inequality fails at
every $n \in \{4,5,\dots,12\}$.
\end{proof}

\begin{corollary}
The conjunction forming conjecture \texttt{00000000443} is false.
\end{corollary}

\begin{proof}
The conjecture asserts a universally quantified inequality; Theorem~\ref{thm:Afails}
exhibits $n = 4$ for which it fails. A conjunction containing a false conjunct is
false.
\end{proof}

\section{Failure (B): the prime-gap formula is false}

The conjecture further asserts that for prime $n$ the gap between the bound and the
exact value is \emph{exactly} half of $2^{(n-1)/2}$, i.e.\ (taking the magnitude of
the gap)
\[
\bigl|B(n) - \dim L_n\bigr| \;=\; \tfrac12 \cdot 2^{(n-1)/2}
\qquad (n \ \text{prime}).
\]
This is a second, logically independent assertion. Table~\ref{tab:gap} evaluates
both sides for the primes $n = 2,3,5,7,11$.

\begin{table}[h]
\centering
\begin{tabular}{r|r|r|r|r|c}
$n$ (prime) & $\dim L_n$ & $B(n)$ & $|B(n)-\dim L_n|$ &
$\tfrac12 2^{(n-1)/2}$ & claim holds? \\
\hline
2  & 1   & 0   & 1   & (not an integer) & ill-defined \\
3  & 2   & 0   & 2   & 1   & \textbf{no} \\
5  & 6   & 8   & 2   & 2   & yes (coincidence) \\
7  & 18  & 48  & 30  & 4   & \textbf{no} \\
11 & 186 & 960 & 774 & 16  & \textbf{no} \\
\end{tabular}
\caption{The prime-gap assertion. It fails at $n = 3, 7, 11$, holds
coincidentally at $n = 5$, and is not even an integer at $n = 2$.}
\label{tab:gap}
\end{table}

\begin{theorem}[failure of the prime-gap formula]
\label{thm:Bfails}
The assertion ``for prime $n$ the gap is exactly half of $2^{(n-1)/2}$'' is false.
Explicitly:
\[
n=3:\quad |B(3)-\dim L_3| = |0-2| = 2 \ne 1 = \tfrac12 2^{1},
\]
\[
n=7:\quad |B(7)-\dim L_7| = |48-18| = 30 \ne 4 = \tfrac12 2^{3},
\]
\[
n=11:\quad |B(11)-\dim L_{11}| = |960-186| = 774 \ne 16 = \tfrac12 2^{5}.
\]
\end{theorem}

\begin{proof}
The values $\dim L_3 = 2$, $\dim L_7 = 18$, $\dim L_{11} = 186$ are read from
Table~\ref{tab:dim}; the bounds are
$B(3) = 0$, $B(7) = 48$, $B(11) = 960$ from Table~\ref{tab:bound}. Each displayed
inequality is then integer arithmetic. Since the claimed identity fails at three
primes, it is false as a general statement.
\end{proof}

\begin{remark}
At $n = 5$ we have $|B(5)-\dim L_5| = |8-6| = 2 = \tfrac12 2^{2}$, so the formula
happens to be correct there; a single coincidence cannot support a claim quantified
over all primes, and $n = 3, 7, 11$ already refute it. At $n = 2$ the quantity
$\tfrac12 2^{(n-1)/2} = \tfrac12\sqrt2$ is irrational, while the gap is the integer
$1$, so the phrase ``exactly half of $2^{(n-1)/2}$'' cannot hold at $n = 2$ under
any reading either.
\end{remark}

\section{The two failures are independent}

Failure (A) is a statement about the \emph{inequality} and is witnessed at the
composite value $n = 4$ (and all $n \ge 4$ in the tested range). Failure (B) is a
statement about the \emph{prime gap} and is witnessed at the primes $n = 3,7,11$.
Either failure alone makes the conjecture false: the conjecture is the conjunction
of (i) the universal inequality and (ii) the prime-gap identity, and a conjunction
is false as soon as one conjunct is false. Accordingly, even a reader who doubted
one of the two computations would still be forced to reject the conjecture by the
other.

\section{Reproducing the computation}

The numerical content is reproduced by \texttt{reproduce.py}, which uses only the
Python~3 standard library and implements the Möbius function by hand:

\begin{verbatim}
python3 reproduce.py
\end{verbatim}

It prints Tables~\ref{tab:dim}--\ref{tab:gap}, checks that
$\dim L_n < B(n)$ for all $n \in [4,12]$, checks that the gap claim fails at
$n = 3,7,11$ (and notes the coincidence at $n = 5$), and exits with status $0$
(\texttt{PASS}) only if every check succeeds.

\end{document}
